% =============================================================
% 4.1 การวิเคราะห์ระบบ
% โครงงาน: การออกแบบและพัฒนาเกมเดินเรือด้วยระบบสั่งการด้วยเสียงเพื่อส่งเสริมทักษะการสื่อสาร
% =============================================================

\section{การวิเคราะห์ระบบ}

\subsection{ผู้ใช้งานระบบและสิทธิ์การเข้าถึง}

เกมเป็นเกมผู้เล่นเดี่ยวที่ทำงานบนเครื่องคอมพิวเตอร์ของผู้เล่น จึงไม่มีระบบบัญชีผู้ใช้หรือการเข้าสู่ระบบ อย่างไรก็ตาม เกมมีผู้เกี่ยวข้อง 2 บทบาท ซึ่งใช้งานส่วนต่าง ๆ ของระบบแตกต่างกัน ดังตารางที่ \ref{tab:roles}

\begin{table}[!htbp]
\centering
\caption{บทบาทผู้ใช้งานและส่วนของระบบที่เข้าถึง}
\label{tab:roles}
\begin{tabular}{|L{3.00cm}|L{10.70cm}|}
\hline
\textbf{บทบาท} & \textbf{ส่วนของระบบที่ใช้งาน} \\
\hline
ผู้เล่น (กัปตัน) & สั่งการลูกเรือด้วยเสียงหรือพิมพ์คำสั่ง สนทนากับลูกเรือ เล่นบทฝึกสอน เลือกงานและอัปเกรดเรือที่อู่ต่อเรือ เล่นโหมดเนื้อเรื่อง ดูหน้าสรุปผล และปรับการตั้งค่าไมโครโฟน เสียง และปุ่มควบคุม \\ \hline
ผู้ดำเนินการทดลอง & ผู้ที่ดำเนินการทดลองกับผู้เข้าร่วมแต่ละคน (ในโครงงานนี้คือผู้จัดทำ) เปิดเกมใหม่สำหรับผู้เข้าร่วมแต่ละคน เปิดแผงแสดงเวลาตอบสนอง (ปุ่ม F8) ระหว่างการทดลอง และนำแฟ้มบันทึกผล CSV จากเครื่องของผู้เล่นไปวิเคราะห์ \\
\hline
\end{tabular}
\end{table}

จากตารางที่ \ref{tab:roles} ผู้ดำเนินการทดลองไม่มีสิทธิ์พิเศษภายในเกม แต่ใช้ส่วนที่ออกแบบไว้สำหรับเก็บข้อมูลการทดลอง ซึ่งไม่รบกวนการเล่นของผู้เล่น

\subsection{ความต้องการเชิงฟังก์ชัน}

ความต้องการเชิงฟังก์ชันของเกมจำแนกตามโมดูล แสดงดังตารางที่ \ref{tab:func_req}--\ref{tab:func_req2}

\begin{table}[!htbp]
\centering
\caption{ความต้องการเชิงฟังก์ชัน (FR-01 ถึง FR-10)}
\label{tab:func_req}
\begin{tabular}{|L{1.18cm}|L{1.76cm}|L{8.33cm}|L{1.57cm}|}
\hline
\textbf{รหัส} & \textbf{โมดูล} & \textbf{รายละเอียด} & \textbf{ผู้ใช้} \\
\hline
FR-01 & วิทยุ & รับเสียงพูดผ่านปุ่มวิทยุ 2 ปุ่ม ได้แก่ ปุ่มสั่งการ (T) และปุ่มสนทนา (Y) แบบกดค้างหรือกดสลับ & ผู้เล่น \\ \hline
FR-02 & วิทยุ & แปลงเสียงพูดภาษาอังกฤษเป็นข้อความด้วยโมเดล Whisper ภายในเครื่อง & ผู้เล่น \\ \hline
FR-03 & วิทยุ & ตรวจสอบสัญญาณเสียงก่อนแปลง (ไม่มีเสียง สั้นเกินไป เบาเกินไป) และกรองข้อความที่เป็นเสียงรบกวน & ผู้เล่น \\ \hline
FR-04 & คำสั่ง & ตัวแยกคำสั่ง (Command Parser) ภายในเครื่องแยกคำสั่งการเดินเรือมาตรฐาน (เครื่องยนต์ หางเสือ คำถาม ยกเลิกคำสั่ง เทียบท่า) ได้สูงสุด 3 คำสั่งต่อประโยค & ผู้เล่น \\ \hline
FR-05 & คำสั่ง & จำแนกเจตนาด้วย GPT-4o mini เมื่อตัวแยกคำสั่งไม่พบคำสั่งที่ชัดเจน & ผู้เล่น \\ \hline
FR-06 & ลูกเรือ & ลูกเรือทวนคำสั่ง ขอให้พูดซ้ำเมื่อไม่เข้าใจ ให้คำใบ้หลังไม่เข้าใจติดกัน 2 ครั้ง และทวนคำสั่งด้วยรูปประโยคมาตรฐาน & ผู้เล่น \\ \hline
FR-07 & ลูกเรือ & สนทนากับลูกเรือผ่าน GPT-4o mini โดยส่งสถานะเรือจริงเป็นบริบทและสั่งให้ตอบจากข้อมูลนั้น & ผู้เล่น \\ \hline
FR-08 & ลูกเรือ & แยกช่องวิทยุของต้นกลเรือ (Harris, ช่องวิทยุ 1) และพลถือท้ายและยามระวังหน้า (Finn, ช่องวิทยุ 2) และแจ้งเมื่อสั่งผิดช่อง & ผู้เล่น \\ \hline
FR-09 & เรือ & จำลองการเคลื่อนที่ของเรือด้วยแรงเฉื่อย อัตราการหันที่เพิ่มขึ้นทีละน้อย (Yaw Acceleration) การไถลด้านข้าง (Lateral Drift) และคันโทรเลขเครื่องยนต์ 8 ระดับ & ผู้เล่น \\ \hline
FR-10 & ด่าน & สุ่มวางสิ่งกีดขวาง (หิน แนวหินโสโครก ทุ่นระเบิด สันดอน) บนตาราง ตรวจด้วยอัลกอริทึม A* ว่ามีเส้นทางจากท่าต้นทางถึงท่าปลายทาง และสุ่มใหม่เมื่อไม่มีเส้นทางหรือเส้นทางตรงเกินไป โดยช่องที่เส้นทางผ่านกว้างไม่น้อยกว่าที่กำหนด & ผู้เล่น \\
\hline
\end{tabular}
\end{table}

\begin{table}[!htbp]
\centering
\caption{ความต้องการเชิงฟังก์ชัน (FR-11 ถึง FR-20)}
\label{tab:func_req2}
\begin{tabular}{|L{1.18cm}|L{1.76cm}|L{8.33cm}|L{1.57cm}|}
\hline
\textbf{รหัส} & \textbf{โมดูล} & \textbf{รายละเอียด} & \textbf{ผู้ใช้} \\
\hline
FR-11 & เตือนภัย & แสดงสิ่งกีดขวางบนเรดาร์ เตือนภัย 3 ระดับตามเวลาก่อนชน และสถานการณ์คำสั่งฉุกเฉินในโหมดเนื้อเรื่อง & ผู้เล่น \\ \hline
FR-12 & การสอน & บทฝึกสอน 9 ขั้นตอนสำหรับการเดินเรือครั้งแรก โดยเรือไม่เสียหายระหว่างฝึก & ผู้เล่น \\ \hline
FR-13 & การสอน & คำแนะนำประโยคตามสถานการณ์ (Try Saying) การ์ดวิทยุ (F1) การอธิบายความหมายคำศัพท์ และรายงานภาษาท้ายรอบ & ผู้เล่น \\ \hline
FR-14 & อู่ต่อเรือ & เลือกงานขนส่ง 3 ระดับ อัปเกรดเรือ 5 ด้าน (เครื่องยนต์ หางเสือ ยามระวังหน้า ตัวเรือ ลูกเรือ) ด้วยเงินและใบรับรองทักษะภาษา และซ่อมเรือโดยไม่มีค่าใช้จ่าย & ผู้เล่น \\ \hline
FR-15 & เนื้อเรื่อง & โหมดเนื้อเรื่อง 3 บท บทละ 3 ภารกิจ พร้อมแผนที่ บทสนทนา และคำศัพท์ประจำภารกิจ & ผู้เล่น \\ \hline
FR-16 & ท่าเรือ & ขอเทียบท่า (Request Docking) และผูกเรือ (Make Fast) เมื่อถึงจุดหมาย & ผู้เล่น \\ \hline
FR-17 & สรุปผล & แสดงผลการเดินเรือและสรุปการสื่อสารภาษาอังกฤษทางวิทยุหลังจบแต่ละรอบ & ผู้เล่น \\ \hline
FR-18 & การวิจัย & บันทึกผลการเล่นทุกคำสั่ง ทุกรอบ และทุกสถานการณ์ฉุกเฉินลงแฟ้ม CSV แยกตามรอบการเปิดเกม และแผงแสดงเวลาตอบสนอง (F8) & ผู้ดำเนินการทดลอง \\ \hline
FR-19 & ทางเลือก & พิมพ์คำสั่งแทนการพูด (ปุ่ม Enter) โดยผ่านขั้นตอนเดียวกับคำสั่งเสียง & ผู้เล่น \\ \hline
FR-20 & บันทึก & บันทึกความก้าวหน้า (เงิน อัปเกรด ใบรับรอง เนื้อเรื่อง) และการตั้งค่าไว้ในเครื่อง & ผู้เล่น \\
\hline
\end{tabular}
\end{table}

จากตารางที่ \ref{tab:func_req}--\ref{tab:func_req2} เกมมีความต้องการเชิงฟังก์ชันทั้งสิ้น 20 ข้อ ครอบคลุม 13 โมดูล

\subsection{ความต้องการที่ไม่ใช่เชิงฟังก์ชัน}

ความต้องการที่ไม่ใช่เชิงฟังก์ชันแสดงดังตารางที่ \ref{tab:nfr}

\begin{table}[!htbp]
\centering
\caption{ความต้องการที่ไม่ใช่เชิงฟังก์ชัน}
\label{tab:nfr}
\begin{tabular}{|L{1.50cm}|L{2.60cm}|L{9.19cm}|}
\hline
\textbf{รหัส} & \textbf{ด้าน} & \textbf{เกณฑ์ที่วัดได้} \\
\hline
NFR-01 & ประสิทธิภาพ & สำหรับคำสั่งที่มีความยาวเสียงไม่เกิน 2 วินาที เรือต้องเริ่มทำตามภายใน 3 วินาทีหลังผู้เล่นปล่อยปุ่มวิทยุ บนเครื่องทดสอบในตารางที่ \ref{tab:hardware} \\ \hline
NFR-02 & ความเป็นส่วนตัว & การแปลงเสียงเป็นข้อความทำงานภายในเครื่องทั้งหมด ไม่ส่งแฟ้มเสียงออกนอกเครื่อง และไม่บันทึกแฟ้มเสียงของผู้เล่น \\ \hline
NFR-03 & ความทนทาน & หากบริการ OpenAI ไม่ตอบภายในเวลาที่กำหนด (จำแนกเจตนา 6 วินาที สนทนา 15 วินาที) หรือไม่มีกุญแจ API เกมต้องเล่นต่อได้ โดยลูกเรือขอให้พูดซ้ำแทน \\ \hline
NFR-04 & ความปลอดภัย & ไม่บันทึกกุญแจ API ของ OpenAI ไว้ในแฟ้มโครงการ แต่อ่านจากตัวแปรสภาพแวดล้อม \texttt{OPENAI\_\allowbreak API\_\allowbreak KEY} หรือจากแฟ้มที่ไม่เก็บในระบบควบคุมเวอร์ชัน \\ \hline
NFR-05 & ใช้งานง่าย & การส่งคำสั่งหนึ่งครั้งมีความยาวเสียง 0.35--8 วินาที เมื่อเกิน 8 วินาทีระบบตัดและส่งให้อัตโนมัติ พร้อมแจ้งผู้เล่นให้พูดสั้นลง \\ \hline
NFR-06 & การเข้าถึง & ผู้เล่นที่ไม่มีไมโครโฟนพิมพ์คำสั่งแทนได้ เปลี่ยนปุ่มวิทยุได้ ปรับขนาดคำบรรยายได้ และรองรับจอยเกม (Gamepad) \\ \hline
NFR-07 & ความถูกต้อง & ทุกประโยคที่เกมแนะนำให้ผู้เล่นพูด ต้องเป็นรูปประโยคมาตรฐานที่ตัวแยกคำสั่งเข้าใจได้ (ตรวจด้วยการทดสอบหน่วย) \\
\hline
\end{tabular}
\end{table}

จากตารางที่ \ref{tab:nfr} ทุกข้อกำหนดมีเกณฑ์ที่ตรวจสอบได้ ผลการวัดแสดงในหัวข้อการทดสอบระบบ

\subsection{กฎการตรวจสอบความถูกต้องของข้อมูล}

ข้อมูลนำเข้าหลักของเกมคือเสียงพูดของผู้เล่น ระบบจึงกำหนดกฎการตรวจสอบตั้งแต่สัญญาณเสียงจนถึงคำตอบจากปัญญาประดิษฐ์ ดังตารางที่ \ref{tab:validation}

\begin{table}[!htbp]
\centering
\caption{กฎการตรวจสอบความถูกต้องของข้อมูล}
\label{tab:validation}
\begin{tabular}{|L{3.50cm}|L{10.20cm}|}
\hline
\textbf{รายการข้อมูล} & \textbf{กฎการตรวจสอบ} \\
\hline
ความยาวเสียง & ต้องยาวไม่น้อยกว่า 0.35 วินาที และไม่เกิน 8 วินาที \\ \hline
ระดับเสียง & ค่า RMS ของสัญญาณต้องไม่น้อยกว่า 0.006 มิฉะนั้นแจ้งว่าเบาเกินไป \\ \hline
ข้อความจาก Whisper & ระบบปฏิเสธข้อความว่างหรือข้อความที่เป็นเสียงรบกวน เช่น ``[INAUDIBLE]'' และแจ้งให้พูดใหม่ \\ \hline
คำสั่งในประโยคเดียว & รับได้สูงสุด 3 คำสั่ง และหากมีคำสั่งขัดแย้งกัน (เช่น เดินหน้าและถอยหลังพร้อมกัน) ลูกเรือจะถามกลับแทนการทำตาม \\ \hline
คำตอบการจำแนกเจตนา & ต้องเป็น JSON ที่มีชื่อคำสั่งตรงกับรายการคำสั่งของเกมเท่านั้น หากไม่ตรงถือว่าไม่เข้าใจ \\ \hline
รหัสผู้เข้าร่วม & ตัดช่องว่างหัวท้าย ยาวไม่เกิน 24 ตัวอักษร ค่าว่างหมายถึงไม่เข้าร่วมการทดลอง \\
\hline
\end{tabular}
\end{table}

\subsection{ตารางตรวจสอบย้อนกลับ}

ตารางตรวจสอบย้อนกลับ (Traceability Matrix) แสดงความสัมพันธ์ระหว่างความต้องการเชิงฟังก์ชัน คลาสหลักที่พัฒนา ข้อมูลที่ใช้ และกรณีทดสอบ ดังตารางที่ \ref{tab:traceability}

\begin{table}[!htbp]
\centering
\caption{ตารางตรวจสอบย้อนกลับ}
\label{tab:traceability}
\begin{tabular}{|L{1.20cm}|L{5.40cm}|L{4.60cm}|L{1.40cm}|}
\hline
\textbf{FR} & \textbf{คลาสหลัก} & \textbf{ข้อมูลที่เกี่ยวข้อง} & \textbf{TC} \\
\hline
FR-01 & RadioMic, RadioInputController & การตั้งค่าปุ่มวิทยุ & TC-01 \\ \hline
FR-02 & RadioMic, WhisperManager & โมเดล ggml-base.en & TC-01 \\ \hline
FR-03 & RadioMic, CommandRouter & -- & TC-02 \\ \hline
FR-04 & CommandParser, CommandVocabulary & -- & TC-03 \\ \hline
FR-05 & IntentClassifier & OpenAI API & TC-04 \\ \hline
FR-06 & CommandRouter, CrewVoice, StandardPhrases & -- & TC-05 \\ \hline
FR-07 & OpenAIClient, CrewPersona, SmallTalk & OpenAI API & TC-06 \\ \hline
FR-08 & CommandRouter, CommandParser & -- & TC-07 \\ \hline
FR-09 & ShipMovement, ShipEngine, ShipHelm & -- & TC-08 \\ \hline
FR-10 & LevelGenerator & -- & TC-09 \\ \hline
FR-11 & RadarSystem, CrewAlertSystem, EmergencyCall & emergencies.csv & TC-10 \\ \hline
FR-12 & TutorialManager & tutorial\_done, tutorial\_step & TC-11 \\ \hline
FR-13 & PhraseHints, Glossary, LanguageLog & -- & TC-12 \\ \hline
FR-14 & ShipyardManager, Progression & career.*, progress.* & TC-13 \\ \hline
FR-15 & StoryCatalog, StoryProgress, StoryDirector & story.done.* & TC-14 \\ \hline
FR-16 & DockingPilot & -- & TC-15 \\ \hline
FR-17 & MissionResultsPanel & -- & TC-16 \\ \hline
FR-18 & ResearchLog & turns.csv, voyages.csv, emergencies.csv & TC-17 \\ \hline
FR-19 & TypedOrders & -- & TC-18 \\ \hline
FR-20 & GameSession, SettingsService & PlayerPrefs & TC-19 \\
\hline
\end{tabular}
\end{table}

จากตารางที่ \ref{tab:traceability} ความต้องการเชิงฟังก์ชันทุกข้อตรวจสอบย้อนกลับได้ตั้งแต่ความต้องการ $\rightarrow$ คลาสที่พัฒนา $\rightarrow$ ข้อมูล $\rightarrow$ กรณีทดสอบ โดย FR-01 และ FR-02 ทดสอบร่วมกันใน TC-01 เนื่องจากเป็นขั้นตอนต่อเนื่องของการรับคำสั่งเสียงครั้งเดียว

\subsection{แผนภาพกรณีใช้งาน}

จากความต้องการเชิงฟังก์ชันทั้ง 20 ข้อ สามารถสร้างแผนภาพกรณีใช้งานแบ่งตามบทบาทได้ดังภาพที่ \ref{fig:uc_player_sea}--\ref{fig:uc_researcher}

\MermaidAdjustBoxOpts{max width=\linewidth, max height=0.65\textheight, center}

\begin{figure}[!htbp]
\centering
\begin{mermaid}
graph LR
    Player(("ผู้เล่น (กัปตัน)"))

    subgraph Game["เกมเดินเรือด้วยระบบสั่งการด้วยเสียง"]
        UC1["FR-01, FR-02: สั่งการด้วยเสียง"]
        UC2["FR-19: พิมพ์คำสั่ง"]
        UC3["FR-07: สนทนากับลูกเรือ"]
        UC4["FR-11: ตอบสนองต่อการเตือนภัย"]
        UC5["FR-16: ขอเทียบท่าและผูกเรือ"]
        UC6["FR-13: ดูคำแนะนำและการ์ดวิทยุ"]
    end

    Player --- UC1
    Player --- UC2
    Player --- UC3
    Player --- UC4
    Player --- UC5
    Player --- UC6
\end{mermaid}
\caption{แผนภาพกรณีใช้งานของผู้เล่น: ระหว่างการเดินเรือ}
\label{fig:uc_player_sea}
\end{figure}

\begin{figure}[!htbp]
\centering
\begin{mermaid}
graph LR
    Player(("ผู้เล่น (กัปตัน)"))

    subgraph Game["เกมเดินเรือด้วยระบบสั่งการด้วยเสียง"]
        UC7["FR-12: เล่นบทฝึกสอน"]
        UC8["FR-14: เลือกงานและอัปเกรดเรือ"]
        UC9["FR-15: เล่นโหมดเนื้อเรื่อง"]
        UC10["FR-17: ดูหน้าสรุปผล"]
        UC11["FR-20: ปรับการตั้งค่า"]
    end

    Player --- UC7
    Player --- UC8
    Player --- UC9
    Player --- UC10
    Player --- UC11
\end{mermaid}
\caption{แผนภาพกรณีใช้งานของผู้เล่น: ก่อนและหลังการเดินเรือ}
\label{fig:uc_player_port}
\end{figure}

จากภาพที่ \ref{fig:uc_player_sea}--\ref{fig:uc_player_port} ผู้เล่นมีกรณีใช้งาน 11 กรณี แบ่งเป็นการสื่อสารกับลูกเรือระหว่างเดินเรือ และการจัดการก่อนและหลังการเดินเรือ

\begin{figure}[!htbp]
\centering
\begin{mermaid}
graph LR
    Researcher(("ผู้ดำเนินการทดลอง"))

    subgraph Game["เกมเดินเรือด้วยระบบสั่งการด้วยเสียง"]
        UC13["FR-18: ดูเวลาตอบสนอง (F8)"]
        UC14["FR-18: เก็บแฟ้มบันทึกผล CSV"]
    end

    Researcher --- UC13
    Researcher --- UC14
\end{mermaid}
\caption{แผนภาพกรณีใช้งานของผู้ดำเนินการทดลอง}
\label{fig:uc_researcher}
\end{figure}

จากภาพที่ \ref{fig:uc_researcher} ผู้ดำเนินการทดลองมีกรณีใช้งาน 2 กรณี ซึ่งใช้สำหรับการเก็บข้อมูลการทดลองกับกลุ่มตัวอย่าง

\subsection{แผนภาพลำดับ}

\subsubsection{กระบวนการสั่งการด้วยเสียง}

กระบวนการตั้งแต่ผู้เล่นพูดคำสั่งจนเรือเปลี่ยนทิศทางแสดงดังภาพที่ \ref{fig:seq_order}

\begin{figure}[!htbp]
\centering
\begin{mermaid}
sequenceDiagram
    actor P as Player
    participant M as RadioMic
    participant W as Whisper (on-device)
    participant R as CommandRouter
    participant C as CommandParser
    participant V as CrewVoice
    participant S as Ship

    P->>M: Hold T and say "starboard twenty"
    P->>M: Release T
    M->>M: Check length and volume
    M->>W: GetTextAsync(audio)
    W-->>M: "Starboard twenty."
    M->>R: RouteCommand(text)
    R->>C: Parse(text)
    C-->>R: TurnStarboard, 20 degrees
    R->>S: Execute order
    R->>V: Read back
    V-->>P: "Starboard rudder, aye! Turning right. Over."
\end{mermaid}
\caption{แผนภาพลำดับกระบวนการสั่งการด้วยเสียง}
\label{fig:seq_order}
\end{figure}

จากภาพที่ \ref{fig:seq_order} ระบบส่งคำสั่งที่ตัวแยกคำสั่งเข้าใจไปควบคุมเรือทันทีโดยไม่ต้องเรียกบริการภายนอก จึงตอบสนองได้เร็วและทำงานได้แม้ไม่มีการเชื่อมต่ออินเทอร์เน็ต

\subsubsection{กระบวนการเมื่อคำสั่งไม่ชัดเจน}

เมื่อตัวแยกคำสั่งไม่พบคำสั่งที่ชัดเจน ระบบจะส่งข้อความให้ GPT-4o mini จำแนกเจตนา ดังภาพที่ \ref{fig:seq_unclear}

\begin{figure}[!htbp]
\centering
\begin{mermaid}
sequenceDiagram
    actor P as Player
    participant R as CommandRouter
    participant C as CommandParser
    participant I as IntentClassifier
    participant O as OpenAI API
    participant L as OpenAIClient (crew chat)
    participant V as CrewVoice

    P->>R: "starbird fitten" (from Whisper)
    R->>C: Parse(text)
    C-->>R: No confident order
    R->>I: Classify(text)
    I->>O: POST chat/completions (JSON mode)
    O-->>I: {"action": ...}
    alt action found
        I-->>R: TurnStarboard
        R->>V: Read back with the standard phrase
    else not found
        I-->>R: None
        R->>L: Converse (crew reply)
        L->>V: Crew reply text
    else error / timeout / no key
        R->>V: Say again (hint after 2 misses)
    end
    V-->>P: Crew reply on the radio
\end{mermaid}
\caption{แผนภาพลำดับกระบวนการเมื่อคำสั่งไม่ชัดเจน}
\label{fig:seq_unclear}
\end{figure}

จากภาพที่ \ref{fig:seq_unclear} หาก GPT-4o mini จำแนกได้ว่าเป็นคำสั่งใด ลูกเรือจะทวนคำสั่งนั้นด้วยรูปประโยคมาตรฐาน หากจำแนกไม่ได้ ระบบจะส่งข้อความต่อไปยังการสนทนากับลูกเรือ (GPT-4o mini) เพื่อให้ลูกเรือตอบตามบริบท ลูกเรือจะขอให้พูดซ้ำ (``Say again, Bridge. Over.'') เฉพาะเมื่อใช้การสนทนาไม่ได้ เช่น ไม่มีกุญแจ API บริการขัดข้อง หรือหมดเวลา และจะให้คำใบ้เมื่อไม่เข้าใจติดกัน 2 ครั้ง

\subsubsection{กระบวนการสนทนากับลูกเรือ}

เมื่อผู้เล่นกดปุ่มสนทนา (Y) ระบบส่งข้อความไปสร้างคำตอบของลูกเรือ ดังภาพที่ \ref{fig:seq_chat}

\begin{figure}[!htbp]
\centering
\begin{mermaid}
sequenceDiagram
    actor P as Player
    participant R as CommandRouter
    participant T as ShipTelemetry
    participant L as OpenAIClient
    participant O as OpenAI API
    participant V as CrewVoice

    P->>R: Hold Y: "Finn, is it safe ahead?"
    R->>T: Capture ship status
    T-->>R: Heading, speed, hazard, weather
    R->>L: Converse(crew, text)
    L->>O: POST chat/completions (persona + status)
    O-->>L: Short crew reply
    L->>V: Subtitle and voice
    V-->>P: "Reef two hundred metres ahead, Sir. Over."
\end{mermaid}
\caption{แผนภาพลำดับกระบวนการสนทนากับลูกเรือ}
\label{fig:seq_chat}
\end{figure}

จากภาพที่ \ref{fig:seq_chat} การสนทนาไม่สามารถสั่งเรือได้ ระบบส่งสถานะเรือจริงเป็นบริบทและสั่งให้ลูกเรือตอบจากข้อมูลนั้น เพื่อลดโอกาสที่โมเดลภาษาจะสร้างข้อมูลที่ไม่มีอยู่จริง

\subsection{รายละเอียดกรณีใช้งาน}

\subsubsection{กรณีใช้งาน: สั่งการด้วยเสียง}

รายละเอียดกรณีใช้งานการสั่งการด้วยเสียงแสดงดังตารางที่ \ref{tab:uc_order}

\begin{table}[!htbp]
\centering
\caption{รายละเอียดกรณีใช้งาน: สั่งการด้วยเสียง}
\label{tab:uc_order}
\begin{tabular}{|L{3.50cm}|L{10.20cm}|}
\hline
\textbf{รายการ} & \textbf{รายละเอียด} \\
\hline
ชื่อกรณีใช้งาน & สั่งการด้วยเสียง (FR-01 ถึง FR-06) \\ \hline
ผู้กระทำ & ผู้เล่น \\ \hline
เงื่อนไขก่อนเริ่ม & อยู่ระหว่างการเดินเรือ และมีไมโครโฟนที่ใช้งานได้ \\ \hline
ขั้นตอนหลัก & 1) ผู้เล่นเลือกช่องวิทยุของลูกเรือ 2) กดปุ่ม T ค้างไว้ (หรือกดหนึ่งครั้งในโหมดกดสลับ) แล้วพูดคำสั่ง 3) ปล่อยปุ่ม (หรือกดอีกครั้ง) 4) ระบบตรวจสัญญาณเสียงและแปลงเป็นข้อความ 5) ตัวแยกคำสั่งระบุคำสั่ง 6) เรือทำตามคำสั่ง 7) ลูกเรือทวนคำสั่งทางวิทยุ \\ \hline
ขั้นตอนทางเลือก & 4a) เสียงสั้นหรือเบาเกินไป: แจ้งผู้เล่นและไม่ส่งคำสั่ง \newline 5a) ไม่พบคำสั่งชัดเจน: ส่งให้ GPT-4o mini จำแนกเจตนา \newline 5b) GPT-4o mini ไม่พบคำสั่ง: ส่งข้อความไปยังการสนทนากับลูกเรือ (GPT-4o mini) และลูกเรือตอบตามบริบท \newline 5c) ใช้การสนทนาไม่ได้ (ไม่มีกุญแจ API ขัดข้อง หรือหมดเวลา): ลูกเรือขอให้พูดซ้ำ และให้คำใบ้เมื่อไม่เข้าใจติดกัน 2 ครั้ง \newline 5d) สั่งผิดช่อง: ลูกเรือแจ้งให้สั่งกับลูกเรือที่รับผิดชอบ \\ \hline
เงื่อนไขหลังจบ & เรือเปลี่ยนความเร็วหรือทิศทางตามคำสั่ง และระบบบันทึกผลลงแฟ้ม turns.csv \\
\hline
\end{tabular}
\end{table}

\subsubsection{กรณีใช้งาน: ตอบสนองต่อสถานการณ์ฉุกเฉิน}

รายละเอียดกรณีใช้งานสถานการณ์ฉุกเฉินในโหมดเนื้อเรื่องแสดงดังตารางที่ \ref{tab:uc_emergency}

\begin{table}[!htbp]
\centering
\caption{รายละเอียดกรณีใช้งาน: ตอบสนองต่อสถานการณ์ฉุกเฉิน}
\label{tab:uc_emergency}
\begin{tabular}{|L{3.50cm}|L{10.20cm}|}
\hline
\textbf{รายการ} & \textbf{รายละเอียด} \\
\hline
ชื่อกรณีใช้งาน & ตอบสนองต่อสถานการณ์ฉุกเฉิน (FR-11) \\ \hline
ผู้กระทำ & ผู้เล่น \\ \hline
เงื่อนไขก่อนเริ่ม & อยู่ในภารกิจโหมดเนื้อเรื่อง และเรือกำลังแล่นเข้าหาหิน แนวหินโสโครก หรือทุ่นระเบิด \\ \hline
ขั้นตอนหลัก & 1) เรดาร์ตรวจพบหิน แนวหินโสโครก หรือทุ่นระเบิดที่เรือจะชนภายในประมาณ 4.5 วินาที 2) เวลาในเกมช้าลง และหน้าจอแสดงคำสั่งที่ต้องพูด เช่น ``hard-a-starboard'' 3) ผู้เล่นพูดคำสั่งหลบภายในเวลาที่กำหนด (ครั้งแรก 6 วินาที ครั้งต่อไป 4.5 วินาที) โดยคำสั่งหยุดเครื่องหรือถอยหลังก็นับว่าทันเช่นกัน 4) เมื่อตอบทันเวลา การสัมผัสสิ่งกีดขวางครั้งถัดไปภายใน 6 วินาที (เวลาในเกม) เป็นเพียงการเฉียดโดยตัวเรือไม่เสียความทนทาน และลูกเรือชมเมื่อใช้รูปประโยคมาตรฐาน หรือทวนรูปประโยคมาตรฐานให้เมื่อใช้ภาษาทั่วไป \\ \hline
ขั้นตอนทางเลือก & 3a) สั่งผิดทิศ: ลูกเรือตอบว่า ``Other way!'' พร้อมย้ำคำสั่งที่ถูกต้อง และนาฬิกายังคงนับต่อ \newline 3b) หมดเวลา: ลูกเรือเตือนให้เตรียมรับแรงกระแทก และเรือชนสิ่งกีดขวางตามปกติ \newline 3c) นาฬิกาหยุดนับขณะวิทยุกำลังบันทึกเสียงและขณะแปลงเสียงเป็นข้อความ เพื่อไม่ให้ความหน่วงของระบบนับเป็นความผิดของผู้เล่น \\ \hline
เงื่อนไขหลังจบ & ระบบบันทึกผล เวลาตอบสนอง และจำนวนครั้งที่พูดลงแฟ้ม emergencies.csv \\
\hline
\end{tabular}
\end{table}

\subsection{แผนภาพสถานะ}

พฤติกรรมของเกมอธิบายด้วยแผนภาพสถานะ 2 ภาพ ภาพแรกเป็นแผนภาพสถานะของระบบเกมหลัก ซึ่งแสดงการเปลี่ยนสถานะจากเมนูหลัก อู่ต่อเรือ การเดินเรือ จนถึงหน้าสรุปผล ดังภาพที่ \ref{fig:state_main} ในภาคผนวก ค

ภาพที่สองเป็นแผนภาพสถานะของการประมวลผลคำสั่งและการโต้ตอบกับลูกเรือ ซึ่งแสดงเส้นทางของเสียงพูดตั้งแต่การบันทึกจนถึงการทวนคำสั่ง การสนทนา หรือการขอให้พูดซ้ำ ดังภาพที่ \ref{fig:state_ai_communication} ในภาคผนวก ค
